% =========================================================
\section{Introduction}
\label{sec:introduction}
% =========================================================

Learning predictive models for nonlinear controlled systems is a central problem in system identification and model-based control \cite{ljung1999system, vanoverschee1994n4sid, nelles2020nonlinear}. For controller design, model quality is determined by prediction accuracy and by compatibility with control synthesis; models that are both accurate and structurally suited to control design are therefore of primary interest.

The Koopman framework addresses this by representing nonlinear dynamics as a linear evolution in a lifted space of observable functions \cite{koopman1931hamiltonian, brunton2022modern}. This linear structure enables standard linear prediction and controller design tools \cite{otto2021koopman, korda2018linear}, making Koopman modeling an attractive paradigm for control. Its effectiveness depends critically on the choice of lifting map.

Most existing approaches, including EDMD \cite{williams2015data} and learned lifting methods \cite{lusch2018deep, takeishi2017koopman}, construct observables from instantaneous states and identify a single time-invariant operator, neglecting short-term temporal correlations. Delay-embedding variants enrich the observables with time-shifted measurements \cite{kamb2020timedelay, brunton2017chaos}, and context-aware formulations regenerate the operator from recent data \cite{liu2023koopa, li2025mamko}; in neither case is the history placed inside a single low-rank transition family identified jointly with the lifting map. Consequently, prediction errors accumulate during multi-step autoregressive rollout.

We propose Transformer-Koopman, a temporal lifting framework with two components. The first is a lifting map implemented as a temporal encoder: it processes a finite history window through a Transformer encoder with non-overlapping multivariate patches and combines the resulting history features with the current physical state to form a structured lifted state, from which the physical state is recovered by a simple linear projection. The second identifies the Koopman dynamics from the lifted state: rather than fitting a single time-invariant operator, we let the history modulate the lifted transition through a low-rank term, so that the same instantaneous state can evolve differently depending on the recent trajectory. We further show that a history-conditioned operator trained only on encoder-consistent one-step transitions develops a training--deployment inconsistency under recursive scheduled propagation---predicted latent states alter the predicted history, which in turn reschedules the subsequent transition operators---and that spectral stabilization alone does not repair this failure mode. We therefore train the predictor under the same recursively scheduled rollout used at deployment, which restores long-horizon accuracy.

Building on the rollout-trained scheduled predictor, we develop a real-time predictive controller based on a one-iteration RTI-MPC that avoids re-running the full encoder at every prediction stage, and we establish a finite-horizon theory whose assumptions are assessed by empirical diagnostics on the validation distribution.

The contributions of this paper are summarized as follows:
\begin{enumerate}
\item \textbf{Low-rank history-conditioned modulation of the lifted transition.} History is used not only to form the lifted state but also to modulate the effective lifted transition, which is confined to a single low-rank family identified jointly with the lifting map. Unlike context-aware or time-varying Koopman formulations that regenerate the operator from data \cite{liu2023koopa, li2025mamko}, and unlike low-rank Koopman operators that replace the transition itself by a rank-$r$ factorization \cite{forootani2026spectral}, the nominal transition remains full-rank and only its history-dependent deviation is low-rank, so linearity in the lifted state is preserved and the predictive-control problem remains a quadratic program. Ablations show that the selective history dependence is concentrated in the transition modulation.
\item \textbf{Scheduled-rollout inconsistency and deployment-consistent training.} We identify a training--deployment inconsistency that arises when a one-step-trained history-conditioned Koopman model is recursively propagated with its context rescheduled from predicted histories: the latent and context distributions progressively depart from those encountered during one-step fitting, and imposing spectral stability does not repair this failure mode. We therefore train the predictor under the same recursively scheduled propagation used at deployment, exposing it to its self-generated latent states, predicted histories, and context sequence. This restores long-horizon scheduled-propagation accuracy to a level close to recurrent re-encoding.
\item \textbf{Real-time predictive deployment.} The rollout-trained scheduled predictor is integrated into a one-iteration RTI-MPC. By avoiding a full encoder evaluation at every prediction stage, the controller runs in real time on the soft-robot platform and tracks a matched exact scheduled NMPC oracle within $1\%$. A finite-horizon theory---finite-history Markovization, latent rollout error propagation, a bounded prediction-error theorem, and a local RTI approximation result---is established and supported by empirical diagnostics.
\end{enumerate}

